% Prüfsatz je Einheit: xelatex -jobname=check_e1 "\def\EIN{e1}% Prüfsatz je Einheit: xelatex -jobname=check_e1 "\def\EIN{e1}% Prüfsatz je Einheit: xelatex -jobname=check_e1 "\def\EIN{e1}% Prüfsatz je Einheit: xelatex -jobname=check_e1 "\def\EIN{e1}\input{check.tex}"
\documentclass[11pt]{article}
\usepackage{mathblatt}
\usepackage[hidelinks]{hyperref}
% Bausteine, die die Vorlage nicht führt (wie in den Rahmendateien)
\newcommand{\zweigkopf}[3]{\par\noindent\hypertarget{#3}{}\einheitenkopf*{#1}\vspace{-1mm}\noindent{\small #2\par}\medskip}
\newcommand{\flaechenbild}[8]{\begin{tikzpicture}[thick]
  \draw (0,0) rectangle (#1+#2,-#3-#4);
  \draw (#1,0) -- (#1,-#3-#4);
  \draw (0,-#3) -- (#1+#2,-#3);
  \node[above] at (#1/2,0) {$#5$};
  \node[above] at (#1+#2/2,0) {$#6$};
  \node[left] at (0,-#3/2) {$#7$};
  \node[left] at (0,-#3-#4/2) {$#8$};
\end{tikzpicture}}
\begin{document}
\blattkopf{Terme und binomische Formeln}{Prüfsatz \EIN}
\weit
\input{\EIN_a}
\clearpage
\input{\EIN_l}
\end{document}
"
\documentclass[11pt]{article}
\usepackage{mathblatt}
\usepackage[hidelinks]{hyperref}
% Bausteine, die die Vorlage nicht führt (wie in den Rahmendateien)
\newcommand{\zweigkopf}[3]{\par\noindent\hypertarget{#3}{}\einheitenkopf*{#1}\vspace{-1mm}\noindent{\small #2\par}\medskip}
\newcommand{\flaechenbild}[8]{\begin{tikzpicture}[thick]
  \draw (0,0) rectangle (#1+#2,-#3-#4);
  \draw (#1,0) -- (#1,-#3-#4);
  \draw (0,-#3) -- (#1+#2,-#3);
  \node[above] at (#1/2,0) {$#5$};
  \node[above] at (#1+#2/2,0) {$#6$};
  \node[left] at (0,-#3/2) {$#7$};
  \node[left] at (0,-#3-#4/2) {$#8$};
\end{tikzpicture}}
\begin{document}
\blattkopf{Terme und binomische Formeln}{Prüfsatz \EIN}
\weit
\input{\EIN_a}
\clearpage
\input{\EIN_l}
\end{document}
"
\documentclass[11pt]{article}
\usepackage{mathblatt}
\usepackage[hidelinks]{hyperref}
% Bausteine, die die Vorlage nicht führt (wie in den Rahmendateien)
\newcommand{\zweigkopf}[3]{\par\noindent\hypertarget{#3}{}\einheitenkopf*{#1}\vspace{-1mm}\noindent{\small #2\par}\medskip}
\newcommand{\flaechenbild}[8]{\begin{tikzpicture}[thick]
  \draw (0,0) rectangle (#1+#2,-#3-#4);
  \draw (#1,0) -- (#1,-#3-#4);
  \draw (0,-#3) -- (#1+#2,-#3);
  \node[above] at (#1/2,0) {$#5$};
  \node[above] at (#1+#2/2,0) {$#6$};
  \node[left] at (0,-#3/2) {$#7$};
  \node[left] at (0,-#3-#4/2) {$#8$};
\end{tikzpicture}}
\begin{document}
\blattkopf{Terme und binomische Formeln}{Prüfsatz \EIN}
\weit
\input{\EIN_a}
\clearpage
\input{\EIN_l}
\end{document}
"
\documentclass[11pt]{article}
\usepackage{mathblatt}
\usepackage[hidelinks]{hyperref}
% Bausteine, die die Vorlage nicht führt (wie in den Rahmendateien)
\newcommand{\zweigkopf}[3]{\par\noindent\hypertarget{#3}{}\einheitenkopf*{#1}\vspace{-1mm}\noindent{\small #2\par}\medskip}
\newcommand{\flaechenbild}[8]{\begin{tikzpicture}[thick]
  \draw (0,0) rectangle (#1+#2,-#3-#4);
  \draw (#1,0) -- (#1,-#3-#4);
  \draw (0,-#3) -- (#1+#2,-#3);
  \node[above] at (#1/2,0) {$#5$};
  \node[above] at (#1+#2/2,0) {$#6$};
  \node[left] at (0,-#3/2) {$#7$};
  \node[left] at (0,-#3-#4/2) {$#8$};
\end{tikzpicture}}
\begin{document}
\blattkopf{Terme und binomische Formeln}{Prüfsatz \EIN}
\weit
\input{\EIN_a}
\clearpage
\input{\EIN_l}
\end{document}
